\chapter{Congruencias y aplicaciones}
\label{md:congruencias}

\section{Aritmética modular}
\label{md:aritmetica-modular}

Para un módulo $m>0$,
\[
 a\equiv b\pmod m\ \eqv\ m\mid(a-b)
 \ \eqv\ \MOD(a,m)=\MOD(b,m).
\]
Es una relación de equivalencia. Sus $m$ clases se representan de
forma única por $0,1,\ldots,m-1$, los residuos mínimos no negativos.
Un \textbf{sistema completo de residuos} contiene exactamente un
representante de cada clase.

Las congruencias con el mismo módulo se pueden sumar, restar y multiplicar;
también se conservan al elevar a potencias enteras no negativas.
La cancelación de un factor exige una condición:
\[
 ha\equiv hb\pmod m,\quad\mcd(h,m)=1
 \quad\imp\quad a\equiv b\pmod m.
\]
Si $d>0$ divide a $m$, una congruencia módulo $m$ implica la misma
congruencia módulo $d$. Para varios módulos positivos,
congruencia en todos ellos equivale a congruencia módulo su mínimo
común múltiplo.

\section{Congruencias lineales e inversos}
\label{md:congruencia-lineal}

La congruencia $ax\equiv b\pmod m$ equivale a la ecuación
diofántica $ax-my=b$. Si $d=\mcd(a,m)$, tiene solución
si y solo si $d\mid b$. En ese caso tiene exactamente $d$ clases
de soluciones módulo $m$:
\[
 x\equiv x_0+k\frac md\pmod m,\qquad k=0,\ldots,d-1.
\]
Para obtener $x_0$ se divide por $d$ y se resuelve
$(a/d)x\equiv b/d\pmod{m/d}$.

Un \textbf{inverso de $a$ módulo $m$} es un entero $u$ con
$au\equiv1\pmod m$. Existe si y solo si $\mcd(a,m)=1$,
y es único como clase. Una identidad $au+mv=1$ lo calcula.
Así, cuando $a$ es coprimo con $m$, la solución es $x\equiv ub\pmod m$.

\begin{ejemplo}
$6x\equiv4\pmod{10}$ se reduce a $3x\equiv2\pmod5$.
El inverso de $3$ es $2$, luego $x\equiv4\pmod5$.
Las dos clases módulo $10$ son $4$ y $9$.
\end{ejemplo}

\section{Teorema chino de los restos}
\label{md:chino}

Sean $m_1,\ldots,m_r>1$ coprimos \emph{dos a dos} y
$M=m_1\cdots m_r$. Para cualesquiera residuos $r_i$, el sistema
\[
 x\equiv r_i\pmod{m_i}\qquad(i=1,\ldots,r)
\]
tiene una única clase de soluciones módulo $M$.
Para construirla, sea $t_i=M/m_i$ y elíjase $u_i$ tal que
$t_iu_i\equiv1\pmod{m_i}$; entonces
\[
 \boxed{x\equiv\sum_{i=1}^r r_it_iu_i\pmod M.}
\]
Cada sumando tiene el residuo requerido en su módulo y es cero en
los restantes. La unicidad usa que una diferencia divisible por todos
los $m_i$ es divisible por su producto.

Si una ecuación tiene la forma $a_ix\equiv b_i\pmod{m_i}$ y
$\mcd(a_i,m_i)=1$, primero se multiplica por su inverso para normalizarla.
El módulo $1$ no impone ninguna restricción y se puede eliminar.

La \textbf{aritmética por residuos} representa una clase módulo $M$
mediante sus residuos en los $m_i$. Sumas y productos se realizan
componente a componente y se reconstruyen por el teorema.
El resultado es una clase módulo $M$; para recuperar un entero concreto
hace falta fijar su intervalo, por ejemplo $0\le x<M$.

\section{Función de Euler, Fermat y Wilson}
\label{md:euler-aritmetico}

La \textbf{función de Euler} $\varphi(m)$ cuenta los enteros
$1\le r\le m$ coprimos con $m$; por convenio $\varphi(1)=1$.
Para un primo $p$ y $\alpha\ge1$,
\[
 \varphi(p)=p-1,\qquad
 \varphi(p^\alpha)=p^\alpha-p^{\alpha-1}.
\]
Si $u,v>0$ son coprimos, $\varphi(uv)=\varphi(u)\varphi(v)$.
En general, para $m>1$,
\[
 \varphi(m)=m\prod_{\substack{p\mid m\\p\text{ primo}}}\left(1-\frac1p\right).
\]

El \textbf{teorema de Euler} establece
\[
 \mcd(a,m)=1\quad\imp\quad a^{\varphi(m)}\equiv1\pmod m.
\]
Multiplicar por $a$ permuta las clases coprimas; al comparar sus
productos se cancela un producto también coprimo con $m$.
El \textbf{pequeño teorema de Fermat} es el caso primo:
$p\nmid a$ implica $a^{p-1}\equiv1\pmod p$.

El \textbf{teorema de Wilson} caracteriza los primos:
\[
 n\ge2:\qquad n\text{ primo}\ \eqv\ (n-1)!\equiv-1\pmod n.
\]
Para un primo, cada residuo no nulo se empareja con su inverso;
solo $1$ y $-1$ pueden coincidir con él.
Las condiciones de Fermat y Euler son esenciales: no se cancelan
factores ni se reducen exponentes por estas fórmulas sin coprimalidad.

\section{Generación de residuos y controles}

Un \textbf{generador lineal congruencial} tiene la forma
\[
 x_{n+1}=\MOD(ax_n+c,m),\qquad m\ge2,
\]
con parámetros enteros y una semilla $0\le x_0<m$.
Como hay $m$ estados y cada uno determina el siguiente, la sucesión
acaba entrando en un ciclo de longitud a lo sumo $m$;
puede existir un tramo inicial anterior al ciclo.
Tiene período completo $m$ si y solo si
\begin{itemize}
 \item $\mcd(c,m)=1$;
 \item $a\equiv1\pmod p$ para todo primo $p\mid m$;
 \item si $4\mid m$, también $a\equiv1\pmod4$.
\end{itemize}

Los controles modulares añaden un dato que obliga a cumplir una
congruencia. Un \textbf{bit de paridad} fija la suma de bits módulo $2$:
detecta un número impar de cambios de bits, pero un número par puede
pasar inadvertido. Un control que se cumple no garantiza ausencia de errores.

\section{Cifrado afín y esquema RSA}

Al representar símbolos por residuos módulo $m\ge2$, el
\textbf{cifrado afín} aplica $x\mapsto ax+b\pmod m$.
Para que sea reversible se requiere $\mcd(a,m)=1$; si $u$ es su inverso,
el descifrado es $y\mapsto u(y-b)\pmod m$.
El cifrado de César corresponde a $a=1$.

En el \textbf{esquema RSA} se eligen primos distintos $p,q$,
$N=pq$ y $\varphi(N)=(p-1)(q-1)$.
Se toman enteros positivos $e,d$ con
\[
 \mcd(e,\varphi(N))=1,\qquad ed\equiv1\pmod{\varphi(N)}.
\]
La clave pública es $(N,e)$ y el exponente privado es $d$.
Un mensaje representado por $0\le m<N$ se cifra como
$c\equiv m^e\pmod N$ y se descifra calculando $c^d\pmod N$.
Para mensajes coprimos con $N$, Euler justifica la recuperación:
si $ed=1+k\varphi(N)$,
\[
 c^d\equiv m^{ed}
       =m\bigl(m^{\varphi(N)}\bigr)^k\equiv m\pmod N.
\]
La justificación por Euler requiere precisamente esa coprimalidad.
